% SPDX-FileCopyrightText: Copyright (c) 2020-2026 Yegor Bugayenko
% SPDX-License-Identifier: MIT

\section{Features}

We selected the most complex features of the most popular
  object-oriented programming languages---C++, Java, C\#, Python, Ruby,
  and JavaScript---and demonstrated how each of them may be represented
  in \eolang{}~\citep{bugayenko2021eolang} objects%
    \footnote{%
    \LaTeX{} sources of this paper are maintained in
    \href{https://github.com/REPOSITORY}{REPOSITORY} GitHub repository,
  the rendered version is
    \href{https://github.com/REPOSITORY/releases/tag/0.0.0}{\ff{0.0.0}}.}.
Where a feature is best illustrated by another language, such as C,
  Kotlin, or PHP, we use that language instead.

Other features---such as operators, loops, generators, procedures,
  variables, code blocks, constants, and selection---are more trivial,
  which is why they are not presented in this paper.

\subsection{Static Methods}
\label{sec:static}

A static method (or static function) is a method defined as a member
  of a class but accessible directly through the class itself,
  rather than through an instance created by its constructor.
For example, this C\# class consists of a single static method.
It may be converted to the following \eolang{} code,
  since a static method is nothing more than a ``global'' function.

\begin{ffcode}
class Utils {
  public static int max(int a, int b) {
    if (a > b) return a;
    return b;
  }
}
(*@ \ffcolumnbreak @*)
if. > [a b] > max
  a.gt b
  a
  b
\end{ffcode}

\subsection{Goto}
\label{sec:goto}

A ``backward'' jump in the C language can be mapped to \eolang{} recursion:

\begin{ffcode}
#include "stdio.h"
void f() {
  int i = 1;
  again:
  i++;
  if (i < 10)
    goto again;
  printf("Finished!");
}
(*@ \ffcolumnbreak @*)
malloc.empty > [] > f
  seq * > [i]
    i.put 1
    again i
    io.stdout "Finished!"
  seq * > [ii] > again
    ii.put (ii.as-number.plus 1)
    if.
      ii.as-number.lt 10
      again ii
      true
\end{ffcode}

A ``forward'' jump in the C language can be mapped to \eolang{}:

\begin{ffcode}
int f(int x) {
  int r = 0;
  if (x == 0)
    goto exit;
  r = 42 / x;
  exit:
  return r;
}
(*@ \ffcolumnbreak @*)
malloc.empty > [x] > f
  seq * > [r] >>
    r.put 0
    if.
      x.eq 0
      true
      r.put (42.div x)
    number r.get
\end{ffcode}

Here, the jump to the \ff{exit} label becomes a conditional that skips
  the division when \ff{x} is zero, so \ff{r} keeps its initial value.

According to the B\"{o}hm–Jacopini theorem~\citep{bohm1966flow},
  any use of the \ff{GOTO} operator may be replaced
  by a combination of
  sequence (the \deff{seq} object),
  selection (the \deff{if} attribute),
  and
  iteration (the \deff{while} object).
Relying on this fact, we may conclude that any C code that uses \ff{GOTO}
  may be represented in \eolang{}.

\subsection{Pointers}
\label{sec:pointers}

This is an example of C code, where the pointer \ff{p} is first incremented
  by seven times \ff{sizeof(book)} (which is equal to 104)
  and then de-referenced to become a struct \ff{book} mapped to memory.
Then, the \ff{title} part of the struct is filled with a string
  and the \ff{price} part is returned as a result of the function \ff{f}.
The algorithm can be represented in \eolang{} on top of the \adeff{malloc} object.

\begin{ffcode}
#include "string.h"
struct book {
  char title[100];
  int price;
};
int f(struct book* p) {
  struct book b = *(p + 7);
  strcpy(b.title, "1984");
  return b.price;
}
(*@ \ffcolumnbreak @*)
[mem offset] > book
  book > [i] > plus
    mem
    offset.plus (i.times 104)
  as-i32. > get-price
    mem.read (offset.plus 100) 4
  mem.write offset s > [s] > set-title
[p] > f
  (book p).plus 7 > b
  seq * > @
    b.set-title "1984"
    b.price
\end{ffcode}

A pointer may not only refer to data in the heap but also
  to executable code in memory (there is no difference between
  ``program'' and ``data'' memory in both x86 and ARM architectures).
This is an example of a C program that calls a function referenced
  by a function pointer, providing two integer arguments to it.
It can easily be mapped to \eolang{} via its partial application:

\begin{ffcode}
typedef int foo(int, int);
int f(foo *p) {
  return (*p)(7, 42);
}
(*@ \ffcolumnbreak @*)
[x y] > foo

[foo] > f
  foo 7 42
\end{ffcode}

It is important to note that the following code
  cannot be represented in \eolang{}:

\begin{ffcode}
int f() {
  int (*p)(int);
  p = (int(*)(int)) 0x761AFE65;
  return (*p) (42);
}
\end{ffcode}

Here, the pointer refers to an arbitrary address in program memory,
  where some executable code is supposed to be located.

This C function, which returns seven (tested with GCC),
  assumes that variables are placed in the stack sequentially
  and uses pointer de-referencing to access them.

\begin{ffcode}
long f() {
  long a = 42;
  long b = 7;
  return *(&a - 1);
}
(*@ \ffcolumnbreak @*)
malloc.empty > [] > f
  seq * > [m]
    m.write 0 7
    m.write 8 42
    number
      m.read (8.minus 8) 8
\end{ffcode}

\subsection{Impure Functions}

A function is pure when, among other qualities,
  it does not have side effects: no mutation of local static variables,
  non-local variables, mutable reference arguments, or input/output streams.
This impure PHP function can be mapped to \eolang{}:

\begin{ffcode}
$a = 42;
function inc() {
  $a = $a + 1;
  return $a;
}
inc(inc($a)); // $a == 44
(*@ \ffcolumnbreak @*)
[a] > inc
  number a.get > ia
  seq * > @
    a.put (ia.plus 1)
    ia
\end{ffcode}

Here, \ff{a} is an abstraction of a memory block.
It must be passed as an argument to the \ff{inc} object.
In PHP, and many other languages, the heap is available in global state,
  while in \eolang{} it must be passed explicitly as an argument.

\subsection{Closures}
\label{sec:closures}

A closure is a function bundled with the variables it captured
  from its defining scope.
This JavaScript function returns a counter that increments
  a captured variable on each call.

\begin{ffcode}
function counter() {
  let c = 0;
  return () => ++c;
}
var next = counter();
next(); // 1
next(); // 2
(*@ \ffcolumnbreak @*)
[c] > counter
  seq * > [] > next
    c.put (c.as-number.plus 1)
    c.as-number
malloc.empty > [] > app
  [c]
    counter c > cnt
    seq * > @
      c.put 0
      cnt.next
      cnt.next
\end{ffcode}

In \eolang{}, a closure is just an object, and the captured variable,
  being mutable, is a \adeff{malloc} block.
The only open question is ownership.
The block cannot belong to the closure \ff{next} itself, nor to the
  \ff{counter} that builds it, because \eolang{} memory is scoped and would
  be released once the maker returns.
It must be owned by an object that outlives the closure---here the \ff{app}
  object---which allocates the block and passes it in.
As with impure functions, \eolang{} grants no implicit access to captured
  state: it must be threaded through explicitly.

\subsection{Classes}
\label{sec:classes}

A class is an extensible program code template for creating objects,
  providing initial values for state (member variables)
  and implementations of behavior (member functions or methods).
The following program contains a Ruby class with a single constructor,
  two methods, and one mutable member variable.
It can be mapped to the following \eolang{} code,
  where a class becomes an object factory.

\begin{ffcode}
class Book
  def initialize(i)
    @id = i
    puts "New book!"
  end
  def path
    "/tmp/${@id}.txt"
  end
  def move(i)
    @id = i
  end
end
(*@ \ffcolumnbreak @*)
[id i] > book
  seq * > [] > initialize!
    id.put i
    io.stdout "New book!"
  seq * > [] > path
    initialize
    tt.sprintf *1
      "/tmp/%s.txt"
      id
  seq * > [i] > move
    initialize
    id.put i
\end{ffcode}

Here, the constructor is represented by the object \ff{initialize},
  whose \ff{!} suffix makes it a constant, so it runs only once
  no matter how many methods call it.
Making an ``instance'' of a book would look like this:

\begin{ffcode}
malloc.empty
  [id]
    book id 42 > b
    seq * > @
      b.move 7
      b.path
\end{ffcode}

\subsection{Encapsulation}
\label{sec:encapsulation}

Encapsulation hides the internal state and helper behavior of an object
  behind a public interface.
In this Java class, the field \ff{rate} and the method \ff{tax} are private,
  while only \ff{total} is exposed to callers.

\begin{ffcode}
class Cart {
  private int rate = 10;
  private int tax(int p) {
    return p / rate; }
  int total(int p) {
    return p + tax(p); } }
(*@ \ffcolumnbreak @*)
[] > cart
  p.plus > [p] > total
    p.div 10
\end{ffcode}

\eolang{} has no \ff{private} or \ff{protected} modifiers: an object either
  names an attribute or it does not.
Here, both \ff{rate} and \ff{tax} dissolve into the anonymous sub-expression
  \ff{p.div 10}, which, having no name, cannot be reached from outside,
  so their hiding is total.
Data hiding is thus achieved by anonymity: whatever is not named cannot
  be addressed.
When a private member must keep a name, such as a mutable field shared
  between methods, its privacy reduces to a naming convention, since
  objects carry no notion of visibility.

\subsection{Static Fields}
\label{sec:staticfields}

A static field is mutable state shared by every instance of a class.
In this Java class, the counter \ff{count} is incremented by the constructor,
  so it holds the number of books ever created.

\begin{ffcode}
class Book {
  static int count = 0;
  Book() {
    count += 1;
  }
}
(*@ \ffcolumnbreak @*)
[count] > book
  count.put > [] > make
    count.as-number.plus 1
malloc.empty > [] > app
  [count]
    seq * > @
      count.put 0
      (book count).make
      (book count).make
      count.as-number
\end{ffcode}

\eolang{} has no class to hold shared state, so the \ff{count} block must be
  owned by an object that outlives every instance: the application object
  \ff{app} itself.
Each \ff{book} receives the same block, and an increment made through one
  instance is visible through all of them.
What other languages hide behind the word \ff{static} becomes explicit---the
  owner of class-level state is simply the top-level object of the program.

\subsection{Destructors}
\label{sec:destructors}

In C++ and some other languages, a destructor is a method
  that is called for a class object when that object passes
  out of scope or is explicitly deleted.
The following code prints both \ff{Alive} and \ff{Dead} texts.
In \eolang{}, the destructor must be called explicitly
  by the owner of the object.

\begin{ffcode}
#include <iostream>
class Foo {
public:
  Foo() { std::cout << "Alive"; }
  ~Foo() { std::cout << "Dead"; };
};
int main() {
  Foo f = Foo();
}
(*@ \ffcolumnbreak @*)
[] > foo
  io.stdout "Alive" > [] > ctor
  io.stdout "Dead" > [] > dtor
[] > main
  foo > f
  seq * > @
    f.ctor
    f.dtor
\end{ffcode}

\subsection{Exceptions}
\label{sec:exceptions}

\broken

Exception handling is the process of responding to the occurrence
  of exceptions: anomalous or exceptional conditions requiring special
  processing---during the execution of a program.
This C++ program utilizes exception handling
  to avoid segmentation fault due to null pointer de-referencing.

\begin{ffcode}
#include <iostream>
int price(int book) {
  if (book == 42) throw "Nope!";
  return 7;
}
void print(int b) {
  try {
    std::cout << price(b);
  } catch (char const* e) {
    std::cout << "Error: " << e;
  }
}
(*@ \ffcolumnbreak @*)
if. > [book cant-return] > price
  book.eq 42
  cant-return "Nope!"
  7
seq * > [b] > print
  io.stdout
    price
      b
      [m]:cant-return
\end{ffcode}
\begin{ffcode}
\end{ffcode}

A Java method may throw a number of either checked or unchecked exceptions:

\begin{ffcode}
void f(int x) throws Exception {
  if (x == 0)
    throw new IOException("foo");
  throw new Exception("bar");
}
void bar() {
  try {
    f(42);
  } catch (IOException e) {
    echo(e.getMessage());
  } catch (Exception e) {
    System.out.exit(1);
  }
}
(*@ \ffcolumnbreak @*)
if. > [x] > f
  x.eq 0
  T "foo"
  T "bar"
\end{ffcode}
\begin{ffcode}
\end{ffcode}

\subsection{Types and Type Casting}
\label{sec:types}

A type system is a logical system comprising a set of rules
  that assign a property called a type to various constructs
  of a computer program, such as variables, expressions, functions, or modules.
The main purpose of a type system is to reduce the possibility
  of bugs in computer programs by defining interfaces
  between different parts of a computer program, and then checking
  that the parts have been connected in a consistent way.
In this example, a Java method \ff{cost} expects an argument
  of type \ff{Book} and compilation would fail if another type is provided.

\begin{ffcode}
int cost(Book b) {
  return 42 + b.price();
}
\end{ffcode}

The restrictions enforced by the Java type system at compile time
  through the \ff{Book} type are enforced in \eolang{}
  by automatic type inference.

\subsection{Reflection}
\label{sec:reflection}

Reflection is the ability of a process to examine, introspect,
  and modify its own structure and behavior.
In the following Python example, the method \ff{hello} of an instance
  of class \ff{Foo} is not called directly on the object but is first
  retrieved through reflection functions \ff{globals} and \ff{getattr},
  and then executed:

\begin{ffcode}
def Foo():
    def hello(self, name):
        print("Hello, %s!" % name)
obj = globals()["Foo"]()
getattr(obj, "hello")("Jeff")
\end{ffcode}

It may be implemented in \eolang{} by creating a global dictionary of
  classes and their methods and filling it up at compile time.
Also, all objects may be equipped with a mechanism of registration in
  the dictionary, which will make them discoverable by reflection.

\subsection{Monkey Patching}

Monkey patching is making changes to a module or class
  while the program is running.
This JavaScript program adds a new method \ff{print} to the object \ff{b}
  after the object has already been instantiated:

\begin{ffcode}
function Book(t) { this.title = t; }
var b = new Book("Object Thinking");
b.print = function() {
  console.log(this.title);
}
b.print();
\end{ffcode}

This program may be translated to \eolang{} by introducing
  a new \ff{Book1} class with the \ff{print} method defined differently.

\subsection{Inheritance}
\label{sec:inheritance}

Inheritance is the mechanism of basing a class upon another class
  while retaining similar implementation.
In this Java code class \ff{Book} inherits all methods
  and non-private attributes from class \ff{Item}.
In \eolang{}, composition may be used instead of inheritance.

\begin{ffcode}
class Item {
  private int p;
  int price() { return p; }
}
class Book extends Item {
  int tax() {
    return price() * 0.1;
  }
}
(*@ \ffcolumnbreak @*)
[p] > item
  p > [] > price
[item] > book
  item.price.times 0.1 > [] > tax
\end{ffcode}

So-called prototype-based programming uses generalized objects,
  which can then be cloned and extended.
For example, this JavaScript program defines two objects,
  where \ff{Item} is the parent object and \ff{Book} is the child
  that inherits the parent through its prototype.

\begin{ffcode}
function Item(p) {
  this.price = p;
}
function Book(p) {
  Item.call(this, p);
  this.tax = function () {
    return this.price * 0.1;
  }
}
var t = new Book(42).tax();
console.log(t); // prints "4.2"
(*@ \ffcolumnbreak @*)
[p] > item
  p > price
[item] > book
  item > @
  ^.price.times 0.1 > [] > tax
(book 42).tax > t
io.stdout (as-fixed t)
\end{ffcode}

Multiple inheritance is a feature of some object-oriented
  programming languages in which an object or class can inherit
  features from more than one parent object or parent class.
In this C++ example, the class \ff{Jack} has both \ff{bark} and
  \ff{listen} methods, inherited from \ff{Dog} and \ff{Friend} respectively.

\begin{ffcode}
#include <iostream>
class Dog {
  virtual void bark() {
    std::cout << "Bark!";
  }
};
class Friend {
  virtual void listen();
};
class Jack: Dog, Friend {
  void listen() override {
    Dog::bark();
    std::cout << "Listen!";
  }
};
(*@ \ffcolumnbreak @*)
[] > dog
  io.stdout "Bark!" > [] > bark
[] > friend
  [] > listen
[] > jack
  dog > d
  friend > f
  seq * > [] > listen
    d.bark
    io.stdout "Listen!"
\end{ffcode}

Here, inherited methods are explicitly listed as attributes
  in the object \ff{jack}.
This is very close to what would happen in the virtual table of
  a class \ff{Jack} in C++.
The \eolang{} object \ff{jack} just makes it explicit.

\subsection{Polymorphism}
\label{sec:polymorphism}

Subtype polymorphism lets a single call site invoke different
  implementations depending on the runtime type of the receiver,
  a mechanism known as dynamic dispatch.
In this Java program, \ff{speak} accepts any \ff{Animal}, yet the sound
  it prints is chosen by the concrete object passed to it.

\begin{ffcode}
class Animal {
  String sound() {
    return "?";
  }
}
class Dog extends Animal {
  String sound() {
    return "Bark";
  }
}
void speak(Animal a) {
  print(a.sound());
}
speak(new Dog());
(*@ \ffcolumnbreak @*)
[] > animal
  "?" > [] > sound
[] > dog
  animal > @
  "Bark" > [] > sound
[a] > speak
  io.stdout a.sound
speak dog
\end{ffcode}

In \eolang{}, a method call is nothing but attribute resolution on the
  receiver, so dispatch needs no special mechanism: \ff{a.sound} always
  finds the \ff{sound} of the object bound to \ff{a}.
The object \ff{dog} decorates \ff{animal} through \ff{@}, yet defines its
  own \ff{sound}, which shadows the inherited one.
No virtual table is needed, because each object carries its own behavior.

\subsection{Abstract Classes and Interfaces}
\label{sec:abstract}

An interface, or an abstract class, declares methods without implementing
  them, leaving the implementation to conforming types.
In this Java example, \ff{Shape} declares \ff{area} and \ff{Square}
  implements it.

\begin{ffcode}
interface Shape {
  int area();
}
class Square implements Shape {
  int side;
  public int area() {
    return side * side;
  }
}
(*@ \ffcolumnbreak @*)
[] > shape
  [] > area
[side] > square
  side.times side > [] > area
\end{ffcode}

In \eolang{}, the interface \ff{shape} is an abstract object whose \ff{area}
  attribute has no body; such an object is incomplete and serves purely
  as a contract.
\ff{square} completes the contract by binding \ff{area}.
No \ff{implements} keyword is required: any object that provides the named
  attributes satisfies the interface, and wherever a \ff{shape} is expected,
  type inference verifies that \ff{area} is bound.
An abstract class is the same idea, with some attributes already
  implemented and others left void.

\subsection{Method Overloading}
\label{sec:overloading}

Method overloading is the ability to create multiple functions with
  the same name but different implementations.
Calls to an overloaded function will run a specific implementation
  of that function appropriate to the context of the call,
  allowing one function call to perform different tasks depending on context.
In this Kotlin program, two functions share the same name,
  and the call \ff{foo(42)} resolves to the one that takes an integer.

\begin{ffcode}
fun foo(a: Int) {}
fun foo(a: Double) {}
foo(42)
(*@ \ffcolumnbreak @*)
[a] > foo-with-int
[b] > foo-with-double
foo-with-int 42
\end{ffcode}

\eolang{} has no overloading, so each variant becomes a distinct object
  with a unique name, and the call selects the one whose name matches
  the argument type.

\subsection{Java Generics}
\label{sec:generics}

Generics extend Java's type system to allow a type or method
  to operate on objects of various types while providing
  compile-time type safety.
For example, this Java class expects another class
  to be specified as \ff{T} before use:

\begin{ffcode}
class Cart<T extends Item> {
  private int total;
  void add(T i) {
    total += i.price();
  }
}
(*@ \ffcolumnbreak @*)
[total] > cart
  number total.get > t
  total.put > [i] > add
    t.plus i.price
\end{ffcode}

The type parameter \ff{T} disappears: \ff{cart} accepts any object that
  understands \ff{price}, while type inference recovers the compile-time
  safety that Java enforced through \ff{T}.

\subsection{C++ Templates}
\label{sec:templates}

Templates are a feature of the C++ programming language
  that lets functions and classes operate with generic types,
  working with many different data types without being rewritten
  for each one.
For example, this function returns the larger of two arguments
  of any comparable type \ff{T}:

\begin{ffcode}
template<typename T> T max(T a, T b) {
  return a > b ? a : b;
}
int x = max(7, 42);
(*@ \ffcolumnbreak @*)
[a b] > max
  if. > @
    a.gt b
    a
    b
max 7 42 > x
\end{ffcode}

The \eolang{} \ff{max} needs no type parameter: the same object works
  for any argument that understands \ff{gt}.

\subsection{Mixins}
\label{sec:mixins}

A mixin is a class that contains methods for use by other classes
  without having to be the parent class of those other classes.
The following code demonstrates how a Ruby module is included in a class:

\begin{ffcode}
module Timing
  def recent?
    @time - Time.now < 24 * 60 * 60
  end
end
def News
  include Timing
  def initialize(t)
    @time = t
  end
end
n = News.new(Time.now)
n.recent?
\end{ffcode}

This code may be represented in \eolang{} by just copying
  the method \ff{recent?} to the object \ff{News} as if it were defined there.

\subsection{Annotations}
\label{sec:annotations}

In Java, an annotation is a form of syntactic metadata
  that can be added to source code; classes, methods, variables, parameters,
  and Java packages may be annotated.
Later, annotations may be retrieved through the Reflection API.
For example, this program analyzes the annotation attached to the class
  of the provided object and changes its behavior depending
  on one of the annotation's attributes.

\begin{ffcode}
interface Item {}
@Ship(true)
class Book implements Item {}
@Ship(false)
class Song implements Item {}
class Cart {
  void add(Item i) {
    /* add the item to the cart */
    if (i.getClass()
      .getAnnotation(Ship.class)
      .value()) {
      /* mark cart as shippable */
    }
  }
}
(*@ \ffcolumnbreak @*)
[] > book
  ship true > a1
[] > song
  ship false > a1
[] > cart
  if. > [i] > add
    i.a1.value
      true
      false
\end{ffcode}

Here, each annotation becomes an attribute of the object: \ff{a1} holds
  the \ff{ship} object, and \ff{cart} reads \ff{i.a1.value} to choose
  its behavior, just as the Java code reads it through reflection.

\subsection{Concurrency}
\label{sec:concurrency}

\broken

Concurrency lets multiple computations run simultaneously, leaving the
  language to coordinate their access to shared memory and their order
  of execution.
This Java program starts two threads that print concurrently,
  so the order of the output is not determined in advance.

\begin{ffcode}
new Thread(() ->
  System.out.print("A")
).start();
new Thread(() ->
  System.out.print("B")
).start();
\end{ffcode}

\eolang{} does not yet define a model for concurrent or parallel execution:
  it has no notion of a thread, a lock, or a scheduler.
Until such a model is designed, programs that rely on nondeterministic
  concurrent execution cannot be represented as \eolang{} objects.

\subsection{Traceability}

A certain amount of semantic information may be lost during the translation
  from a more powerful programming language to \eolang{} objects, such as,
  for example, names, line numbers, comments, etc.
Moreover, it may be useful to have the ability to trace \eolang{} objects
  back to the original language constructs from which they were derived.
For example, a simple C function can be mapped
  to the following \eolang{} object.

\begin{ffcode}
int f(int x) {
  return 42 / x;
}
(*@ \ffcolumnbreak @*)
[x] > f
  [] > @
    "src/main.c:1-1" > source
    42.div x > @
  "src/main.c:0-2" > source
\end{ffcode}

Here, the synthetic \ff{source} attribute represents the location
  in the source code file with the original C code.
It is important to ensure during translation that the name \ff{source}
  does not conflict with a possibly similar name in the object.
